% Options for packages loaded elsewhere
\PassOptionsToPackage{unicode}{hyperref}
\PassOptionsToPackage{hyphens}{url}
\documentclass[
]{article}
\usepackage{xcolor}
\usepackage{amsmath,amssymb}
\setcounter{secnumdepth}{-\maxdimen} % remove section numbering
\usepackage{iftex}
\ifPDFTeX
  \usepackage[T1]{fontenc}
  \usepackage[utf8]{inputenc}
  \usepackage{textcomp} % provide euro and other symbols
\else % if luatex or xetex
  \usepackage{unicode-math} % this also loads fontspec
  \defaultfontfeatures{Scale=MatchLowercase}
  \defaultfontfeatures[\rmfamily]{Ligatures=TeX,Scale=1}
\fi
\usepackage{lmodern}
\ifPDFTeX\else
  % xetex/luatex font selection
\fi
% Use upquote if available, for straight quotes in verbatim environments
\IfFileExists{upquote.sty}{\usepackage{upquote}}{}
\IfFileExists{microtype.sty}{% use microtype if available
  \usepackage[]{microtype}
  \UseMicrotypeSet[protrusion]{basicmath} % disable protrusion for tt fonts
}{}
\makeatletter
\@ifundefined{KOMAClassName}{% if non-KOMA class
  \IfFileExists{parskip.sty}{%
    \usepackage{parskip}
  }{% else
    \setlength{\parindent}{0pt}
    \setlength{\parskip}{6pt plus 2pt minus 1pt}}
}{% if KOMA class
  \KOMAoptions{parskip=half}}
\makeatother
\usepackage{color}
\usepackage{fancyvrb}
\newcommand{\VerbBar}{|}
\newcommand{\VERB}{\Verb[commandchars=\\\{\}]}
\DefineVerbatimEnvironment{Highlighting}{Verbatim}{commandchars=\\\{\}}
% Add ',fontsize=\small' for more characters per line
\newenvironment{Shaded}{}{}
\newcommand{\AlertTok}[1]{\textcolor[rgb]{1.00,0.00,0.00}{\textbf{#1}}}
\newcommand{\AnnotationTok}[1]{\textcolor[rgb]{0.38,0.63,0.69}{\textbf{\textit{#1}}}}
\newcommand{\AttributeTok}[1]{\textcolor[rgb]{0.49,0.56,0.16}{#1}}
\newcommand{\BaseNTok}[1]{\textcolor[rgb]{0.25,0.63,0.44}{#1}}
\newcommand{\BuiltInTok}[1]{\textcolor[rgb]{0.00,0.50,0.00}{#1}}
\newcommand{\CharTok}[1]{\textcolor[rgb]{0.25,0.44,0.63}{#1}}
\newcommand{\CommentTok}[1]{\textcolor[rgb]{0.38,0.63,0.69}{\textit{#1}}}
\newcommand{\CommentVarTok}[1]{\textcolor[rgb]{0.38,0.63,0.69}{\textbf{\textit{#1}}}}
\newcommand{\ConstantTok}[1]{\textcolor[rgb]{0.53,0.00,0.00}{#1}}
\newcommand{\ControlFlowTok}[1]{\textcolor[rgb]{0.00,0.44,0.13}{\textbf{#1}}}
\newcommand{\DataTypeTok}[1]{\textcolor[rgb]{0.56,0.13,0.00}{#1}}
\newcommand{\DecValTok}[1]{\textcolor[rgb]{0.25,0.63,0.44}{#1}}
\newcommand{\DocumentationTok}[1]{\textcolor[rgb]{0.73,0.13,0.13}{\textit{#1}}}
\newcommand{\ErrorTok}[1]{\textcolor[rgb]{1.00,0.00,0.00}{\textbf{#1}}}
\newcommand{\ExtensionTok}[1]{#1}
\newcommand{\FloatTok}[1]{\textcolor[rgb]{0.25,0.63,0.44}{#1}}
\newcommand{\FunctionTok}[1]{\textcolor[rgb]{0.02,0.16,0.49}{#1}}
\newcommand{\ImportTok}[1]{\textcolor[rgb]{0.00,0.50,0.00}{\textbf{#1}}}
\newcommand{\InformationTok}[1]{\textcolor[rgb]{0.38,0.63,0.69}{\textbf{\textit{#1}}}}
\newcommand{\KeywordTok}[1]{\textcolor[rgb]{0.00,0.44,0.13}{\textbf{#1}}}
\newcommand{\NormalTok}[1]{#1}
\newcommand{\OperatorTok}[1]{\textcolor[rgb]{0.40,0.40,0.40}{#1}}
\newcommand{\OtherTok}[1]{\textcolor[rgb]{0.00,0.44,0.13}{#1}}
\newcommand{\PreprocessorTok}[1]{\textcolor[rgb]{0.74,0.48,0.00}{#1}}
\newcommand{\RegionMarkerTok}[1]{#1}
\newcommand{\SpecialCharTok}[1]{\textcolor[rgb]{0.25,0.44,0.63}{#1}}
\newcommand{\SpecialStringTok}[1]{\textcolor[rgb]{0.73,0.40,0.53}{#1}}
\newcommand{\StringTok}[1]{\textcolor[rgb]{0.25,0.44,0.63}{#1}}
\newcommand{\VariableTok}[1]{\textcolor[rgb]{0.10,0.09,0.49}{#1}}
\newcommand{\VerbatimStringTok}[1]{\textcolor[rgb]{0.25,0.44,0.63}{#1}}
\newcommand{\WarningTok}[1]{\textcolor[rgb]{0.38,0.63,0.69}{\textbf{\textit{#1}}}}
\setlength{\emergencystretch}{3em} % prevent overfull lines
\providecommand{\tightlist}{%
  \setlength{\itemsep}{0pt}\setlength{\parskip}{0pt}}
\usepackage{microtype}
\usepackage{titlesec}
\usepackage{booktabs}
\usepackage{array}
\usepackage{longtable}
\usepackage{ragged2e}
\usepackage{etoolbox}
\usepackage{caption}
\usepackage{fancyhdr}
\usepackage{setspace}
\usepackage{newunicodechar}

\newunicodechar{✓}{$\checkmark$}
\newunicodechar{∈}{$\in$}
\newunicodechar{∃}{$\exists$}
\newunicodechar{≠}{$\neq$}
\newunicodechar{≤}{$\le$}
\newunicodechar{≥}{$\ge$}
\newunicodechar{→}{$\rightarrow$}
\newunicodechar{₁}{$_1$}
\newunicodechar{₂}{$_2$}
\newunicodechar{ₖ}{$_k$}
\newunicodechar{─}{-}
\newunicodechar{└}{+}
\newunicodechar{├}{+}
\newunicodechar{│}{|}

\setstretch{1.03}
\setlength{\parindent}{0pt}
\setlength{\parskip}{0.45em}
\setlength{\emergencystretch}{3em}

\titleformat{\section}{\Large\bfseries}{\thesection.}{0.6em}{}
\titleformat{\subsection}{\large\bfseries}{\thesubsection.}{0.6em}{}
\titleformat{\subsubsection}{\normalsize\bfseries}{\thesubsubsection.}{0.6em}{}

\titlespacing*{\section}{0pt}{1.4em}{0.6em}
\titlespacing*{\subsection}{0pt}{1.1em}{0.4em}
\titlespacing*{\subsubsection}{0pt}{0.8em}{0.3em}

\setlength{\LTpre}{0.35em}
\setlength{\LTpost}{0.35em}
\setlength{\LTleft}{0pt}
\setlength{\LTright}{0pt}
\setlength{\tabcolsep}{4pt}
\renewcommand{\arraystretch}{1.08}
\AtBeginEnvironment{longtable}{\small}
\newcolumntype{L}[1]{>{\RaggedRight\arraybackslash}p{#1}}
\newcolumntype{C}[1]{>{\Centering\arraybackslash}p{#1}}
\newcolumntype{R}[1]{>{\RaggedLeft\arraybackslash}p{#1}}

\captionsetup[table]{skip=4pt}
\captionsetup[figure]{skip=4pt}

\pagestyle{fancy}
\fancyhf{}
\fancyfoot[C]{\thepage}
\renewcommand{\headrulewidth}{0pt}
\renewcommand{\footrulewidth}{0pt}
\setlength{\headheight}{14pt}

\newcommand{\codeid}[1]{{\small\begingroup\urlstyle{tt}\nolinkurl{#1}\endgroup}}
\usepackage{bookmark}
\IfFileExists{xurl.sty}{\usepackage{xurl}}{} % add URL line breaks if available
\urlstyle{same}
\hypersetup{
  hidelinks,
  pdfcreator={LaTeX via pandoc}}

\author{}
\date{}

\begin{document}

\section{AppSec Core: A Normalized Ontology for Security Requirements
Across Heterogeneous
Frameworks}\label{appsec-core-a-normalized-ontology-for-security-requirements-across-heterogeneous-frameworks}

\textbf{Pedro Farinha}\\
Independent Researcher\\
\href{mailto:pedro.farinha@shiftleft.pt}{\nolinkurl{pedro.farinha@shiftleft.pt}}\\
\href{https://github.com/pedrofarinhaatshiftleftpt}{github.com/pedrofarinhaatshiftleftpt}

\emph{Supported by Shiftleft - Secure Software Engineering, lda.}

\begin{center}\rule{0.5\linewidth}{0.5pt}\end{center}

\subsection{Abstract}\label{abstract}

The application security (AppSec) landscape is served by an expanding
set of frameworks, standards, and regulations, each using its own
vocabulary, granularity, and structural conventions. Organizations that
must demonstrate alignment with multiple sources simultaneously face a
quadratic mapping problem with no shared semantic model to mediate it.

We propose \textbf{AppSec Core}, a normalized ontology for application
security comprising 10 domain slices and 234 typed instances. The
ontology defines four entity types --- ControlObjective, Practice,
Mechanism, and Artifact --- connected by a stable set of relations
(\texttt{implements}, \texttt{realizes}, \texttt{evidences}) that are
structurally invariant across all slices. Each slice is governed by a
formal contract specifying scope, boundaries, and non-goals. The
ontology was derived from practitioner semantics accumulated over a
decade of AppSec engineering, not from any single external framework.

We present evidence that AppSec Core can function as a \textbf{semantic
normalization layer} through canonical reduction: requirements from a
first-wave set of five sources (SSDF, ASVS, SLSA, CIS Controls, and
CAPEC) --- expressed in incompatible vocabularies --- are reduced to
shared ControlObjectives and their associated operational chains within
the ontology, after which downstream coverage analysis operates on the
normalized objective layer independently of framework vocabulary.
Mapping exercises across three presented slices show that 79\% of
analyzed ControlObjectives are mapped by at least two frameworks, and
the 10-slice model accommodates every requirement analyzed under the
stated inclusion protocol without structural extension.

We contribute: (1) the AppSec Core ontology as a reusable domain model;
(2) the four-type instance schema with explicit traceability relations;
(3) empirical evidence of semantic normalization through canonical
reduction --- heterogeneous framework requirements are reduced to shared
ControlObjectives, after which downstream coverage analysis is
independent of framework vocabulary; and (4) a contract-based governance
model that supports principled extension.

\textbf{Keywords:} ontology, application security, normalization,
security frameworks, SSDF, ASVS, SLSA, knowledge representation

\begin{center}\rule{0.5\linewidth}{0.5pt}\end{center}

\subsection{1. Introduction}\label{introduction}

Organizations developing software in regulated or security-sensitive
environments must align their practices with multiple external
frameworks. NIST's Secure Software Development Framework (SSDF) {[}1{]}
prescribes 20 practices across the secure development lifecycle. OWASP's
Application Security Verification Standard (ASVS) v5.0 {[}2{]} defines
443 verification requirements organized in 36 thematic clusters.
Supply-chain Levels for Software Artifacts (SLSA) {[}3{]} specifies 14
build integrity requirements across three maturity levels. CIS Controls
v8.1 {[}4{]} addresses 18 enterprise security controls. Regulatory
instruments such as the EU's DORA {[}5{]} and NIS2 {[}6{]} add further
obligations; their specific mapping to AppSec Core is outside the
evidentiary scope of this paper (see Scope, below).

Each of these sources uses its own terminology, structural conventions,
and level of granularity. SSDF speaks of \emph{practices} organized in
four families (Prepare, Protect, Produce, Respond). ASVS speaks of
\emph{verification requirements} grouped by thematic cluster. SLSA
speaks of \emph{build levels} with specific provenance and isolation
properties. These are not interchangeable vocabularies: they partition
the same underlying domain differently.

For any single framework, alignment is manageable. The problem emerges
at the intersection: \textbf{when an organization must assess coverage
against multiple frameworks simultaneously}, the mapping effort scales
as N internal practices × M external requirements, terminology conflicts
accumulate, and there is no principled way to determine whether two
requirements from different frameworks address the same underlying
security concern.

\subsubsection{The Missing Layer}\label{the-missing-layer}

What is missing is a \textbf{normalization layer} --- a shared semantic
model into which requirements from heterogeneous frameworks can be
mapped, enabling:

\begin{itemize}
\tightlist
\item
  \textbf{Semantic deduplication}: recognizing that SSDF PW.9 (``secure
  settings by default''), ASVS \texttt{secure\_configuration\_baseline},
  and CIS-4 (``secure configuration'') address the same concern
\item
  \textbf{Unified coverage analysis}: determining whether an
  organization's practices cover a requirement \emph{once}, regardless
  of how many frameworks express it
\item
  \textbf{Framework-version resilience}: when ASVS v6.0 is published,
  the mapping to the normalization layer is updated --- downstream
  analysis does not change
\item
  \textbf{Cross-framework gap detection}: identifying requirements
  covered by one framework but not another, through the shared model
\end{itemize}

This paper proposes \textbf{AppSec Core} as this normalization layer.
AppSec Core is an ontology of 10 domain slices with 234 typed instances,
derived from practitioner semantics rather than from any single external
framework.

\subsubsection{Contributions}\label{contributions}

\begin{enumerate}
\def\labelenumi{\arabic{enumi}.}
\item
  \textbf{The AppSec Core ontology}: 10 slices covering the application
  security domain, each with typed instances (ControlObjective,
  Practice, Mechanism, Artifact) connected by explicit relations.
\item
  \textbf{A four-type instance model} with structural invariance across
  all slices, providing a reusable schema for security knowledge
  representation.
\item
  \textbf{Empirical evidence of semantic normalization through canonical
  reduction}: we show through concrete mapping examples that
  requirements from SSDF, ASVS, and SLSA are reduced to shared AppSec
  Core ControlObjectives, producing a canonical representation in which
  framework-specific vocabulary is no longer required for downstream
  coverage analysis.
\item
  \textbf{A contract-based governance model} for ontology extension,
  where each slice is governed by a formal contract specifying scope,
  non-goals, and boundary conditions.
\end{enumerate}

\subsubsection{Scope}\label{scope}

This paper is about \textbf{representation}. It does not address
retrieval pipelines, LLM grounding, or experimental evaluation --- those
are treated in companion work. The empirical mapping evidence reported
here is restricted to the first-wave framework set described above;
later maturity and regulatory waves are outside the evidentiary scope of
this paper. The SbD-ToE manual, from which the ontology was empirically
derived, is referenced as context for the derivation process but is not
the subject of the paper.

\begin{center}\rule{0.5\linewidth}{0.5pt}\end{center}

\subsection{2. Related Work}\label{related-work}

\subsubsection{2.1 Application Security
Frameworks}\label{application-security-frameworks}

SSDF {[}1{]}, ASVS {[}2{]}, SLSA {[}3{]}, and CIS Controls {[}4{]} each
provide comprehensive guidance within their respective scopes. They are
designed as \textbf{consumption endpoints} --- organizations implement
against them --- not as normalization layers. Each defines its own
vocabulary and structure. Cross-framework alignment is left to the
consumer.

\subsubsection{2.2 Security Ontologies}\label{security-ontologies}

Several ontologies have been proposed for information security. Herzog
et al.~{[}7{]} proposed a general security ontology covering assets,
threats, and countermeasures. Blanco et al.~{[}8{]} developed an
ontology for security requirements engineering and provided a systematic
comparison of security ontologies available at the time. These works
operate at the broad information security level; AppSec Core is scoped
specifically to \textbf{application security practices} --- what
development teams do to build secure software --- a narrower but more
operationally precise domain.

CWE {[}9{]} and CAPEC {[}10{]} model the \textbf{problem space} (what
can go wrong). AppSec Core models the \textbf{solution space} (what
practitioners do to prevent it). This distinction is relevant: coverage
analysis requires a solution-space model. CAPEC is included in the
normalization mapping (Section 5) as a problem-space cross-reference:
its attack patterns validate AppSec Core solution-space objectives
rather than prescribe practices, and it is treated as a complementary
adversarial source rather than a normative framework equivalent to SSDF
or ASVS.

\subsubsection{2.3 Maturity Models}\label{maturity-models}

OWASP SAMM {[}11{]} and BSIMM {[}12{]} assess \textbf{how well} an
organization performs security activities at varying maturity levels.
They are complementary to AppSec Core: SAMM/BSIMM assess maturity;
AppSec Core models \textbf{what the activities are}. An organization
could use AppSec Core for coverage analysis and SAMM for maturity
assessment over the same underlying practices.

\subsubsection{2.4 Requirements
Normalization}\label{requirements-normalization}

The NIST SP 800-53 overlay methodology {[}13{]} provides a mechanism for
customizing security control baselines. Fabian et al.~{[}14{]} compare
security requirements engineering methods including SQUARE and SREP.
Goknil et al.~{[}15{]} formalize trace relation semantics for
requirements models. AppSec Core differs from these in providing a
\textbf{semantic normalization layer} with typed instances and explicit
relations, rather than a methodology or a single framework's
customization mechanism.

\subsubsection{2.5 Gap}\label{gap}

To our knowledge, no existing model provides a typed, structurally
invariant, framework-independent semantic layer at the application
security practice level that supports cross-framework normalization.
AppSec Core is designed to address this gap.

\begin{center}\rule{0.5\linewidth}{0.5pt}\end{center}

\subsection{3. The AppSec Core Ontology}\label{the-appsec-core-ontology}

\subsubsection{3.1 Design Principles}\label{design-principles}

AppSec Core was designed to satisfy four principles:

\textbf{P1 --- Practitioner fidelity.} The ontology reflects how
practitioners organize and reason about application security. It was
derived bottom-up from 10+ years of hands-on AppSec engineering, not
from the categories of any specific framework.

\textbf{P2 --- Framework independence.} The ontology is designed to be
stable as external frameworks evolve. Adding a new framework requires
mapping its requirements to existing slices, not restructuring the
ontology.

\textbf{P3 --- Structural invariance.} Every slice uses the same entity
types, the same relation model, and the same contract structure. This
enables uniform tooling and analysis without per-slice customization.

\textbf{P4 --- Contract-based governance.} Each slice is governed by a
formal contract specifying scope, non-goals, and boundary conditions,
enabling principled extension without scope creep.

\subsubsection{3.2 Derivation Process}\label{derivation-process}

The ontology was derived through an iterative, manual-first, qualitative
normalization process over the SbD-ToE manual corpus. The derivation was
not a top-down design exercise, nor was it driven by unsupervised
clustering or topic modelling. No statistical clustering algorithm was
used to define ontology boundaries or slice membership. Ontology
boundaries, slice contracts, and objective atomicity were established by
expert qualitative analysis of requirements and broad control surfaces
across chapters.

The derivation proceeded as follows:

\begin{enumerate}
\def\labelenumi{\arabic{enumi}.}
\tightlist
\item
  \textbf{Corpus analysis}: 4,139 structural units from the manual were
  indexed with document role and normative weight annotations, providing
  a structured surface for qualitative reading
\item
  \textbf{Concern clustering}: recurring security concerns were
  identified by expert reading of requirements and broad control
  surfaces across chapters, then grouped into candidate subdomains by
  semantic judgment --- not by co-occurrence analysis or embedding-based
  clustering
\item
  \textbf{Objective extraction}: each candidate subdomain was opened
  serially as a slice; within each slice, ControlObjectives were
  formulated atomic-first --- each objective as a single, normatively
  coherent statement before practices or mechanisms were added
\item
  \textbf{Practice/mechanism/artifact identification}: operational,
  technical, and evidentiary instances were extracted from the corpus
  and linked to objectives only after the objective layer was coherent
\item
  \textbf{Contract formalization}: each slice advanced only when it
  reached an explicit contract state specifying scope, non-goals, and
  boundary conditions
\end{enumerate}

This method was chosen because the derivation goal was normative, not
descriptive: AppSec Core is not a statistical reflection of the corpus
but a normalization layer with principled boundaries, atomic objectives,
and explicit contracts. Unsupervised methods capture lexical proximity
and semantic similarity, but do not determine what constitutes an atomic
objective, where a legitimate slice boundary lies, or what should remain
outside the core. These are judgment-dependent decisions that require
interpretive expertise, not statistical inference. Future work may use
embedding-based clustering or topic modelling as a post-hoc
triangulation layer for coherence and coverage verification --- but not
as the primary derivation mechanism.

The manual served as the empirical corpus; the ontology is the
abstraction. The ontology is \textbf{not} a description of the manual
--- it is a domain model that the manual's content instantiates.
Evidence for this independence is provided in Section 7.4, where we show
that several ontology concepts were documented by practitioners before
the corresponding framework requirements were published.

\subsubsection{3.3 The Ten Domain Slices}\label{the-ten-domain-slices}

AppSec Core organizes the application security domain into 10 slices.
Table 1 lists each slice with its scope and instance counts.

\textbf{Table 1.} AppSec Core v0 slices.

\begin{longtable}{@{}>{\RaggedRight\arraybackslash}p{1.15cm}>{\RaggedRight\arraybackslash}p{4.25cm}>{\RaggedRight\arraybackslash}p{5.45cm}cccc@{}}
\toprule
\textbf{ID} & \textbf{Slice Name} & \textbf{Scope} & \textbf{CO} & \textbf{Pr} & \textbf{Me} & \textbf{Ar} \\
\midrule
\endfirsthead
\toprule
\textbf{ID} & \textbf{Slice Name} & \textbf{Scope} & \textbf{CO} & \textbf{Pr} & \textbf{Me} & \textbf{Ar} \\
\midrule
\endhead
\bottomrule
\endfoot
\bottomrule
\endlastfoot
ASC-01 & Supply Chain \& Build Integrity & Dependencies, provenance, build security, artifact signing & 7 & 7 & 5 & 10 \\
ASC-02 & Identity, Access \& Session Trust & Authentication, authorization, session management, credentials & 7 & 6 & 6 & 6 \\
ASC-03 & Architecture \& Trust Boundaries & Secure design, trust zones, boundary controls, design review & 7 & 7 & 5 & 4 \\
ASC-04 & Testing, Security Validation \& Empirical Assurance & Testing strategy, risk-proportional criteria, evidence & 7 & 7 & 5 & 6 \\
ASC-05 & Threat Modeling, Risk Disposition \& Mitigation Traceability & Threat analysis, risk treatment, mitigation tracking & 7 & 7 & 6 & 4 \\
ASC-06 & Secret Handling, Protected Configuration \& Operational Identities & Secrets, secure config baselines, machine identity, drift & 7 & 6 & 4 & 2 \\
ASC-07 & Input Validation, Safe Parsing \& Controlled Failure & Input handling, injection prevention, fail-safe, error management & 7 & 6 & 4 & 6 \\
ASC-08 & Integration Trust \& Service-to-Service Security & API security, service mesh, inter-service authentication, contracts & 7 & 6 & 5 & 5 \\
ASC-09 & Release Promotion, Controlled Rollout \& Rollback Readiness & Release gates, progressive deployment, canary, rollback & 7 & 6 & 4 & 8 \\
ASC-10 & Security Event Logging, Audit Trail \& Centralized Logging & Security logging, audit trails, log integrity, observability & 7 & 5 & 4 & 2 \\
 & \textbf{Total} &  & \textbf{70} & \textbf{63} & \textbf{48} & \textbf{53} \\
\end{longtable}

\emph{CO = ControlObjective, Pr = Practice, Me = Mechanism, Ar =
Artifact. Total: 234 instances.}

Each slice contains exactly 7 control objectives (6 atomic + 1
composite). This regularity is a \textbf{methodological choice} for
uniformity in v0, explicitly declared as non-ontological: future
versions may use different cardinalities if the domain requires it.

\subsubsection{3.4 Entity Types and
Relations}\label{entity-types-and-relations}

The ontology defines four core entity types and one supporting type
(EvidencePattern, not core-driving in v0), connected by a stable set of
relations.

\paragraph{Entity Types}\label{entity-types}

\textbf{ControlObjective.} A reusable domain constraint or assurance
goal. Required fields: \texttt{objective\_id}, \texttt{name},
\texttt{objective\_kind} (atomic \textbar{} composite),
\texttt{objective\_type} (governance \textbar{} preventive \textbar{}
detective \textbar{} corrective), \texttt{domain\_key},
\texttt{statement}, \texttt{expected\_outcome},
\texttt{verification\_posture}. Control objectives are the normative
anchors of the ontology --- the ``what must be achieved.''

\textbf{Practice.} A reusable human or process discipline that realizes
objectives, independent of specific tools. Required fields:
\texttt{practice\_id}, \texttt{name}, \texttt{practice\_family},
\texttt{local\_practice\_type}. Practice families are cross-slice:
\texttt{governance\_and\_review},
\texttt{policy\_and\_gate\_enforcement},
\texttt{identity\_and\_trust\_control},
\texttt{validation\_and\_analysis},
\texttt{integrity\_traceability\_and\_records},
\texttt{rollout\_recovery\_and\_runtime\_records}.

\textbf{Mechanism.} A concrete technical or enforceable means that
implements practices or directly realizes objectives. Required fields:
\texttt{mechanism\_id}, \texttt{name}, \texttt{mechanism\_family},
\texttt{local\_mechanism\_type}.

\textbf{Artifact.} A control-relevant or evidence-bearing output used
for review, traceability, and assurance. Required fields:
\texttt{artifact\_id}, \texttt{name}, \texttt{canonical\_role}.
Canonical roles are cross-slice: \texttt{configuration},
\texttt{governance\_record}, \texttt{review\_record},
\texttt{approval\_record}, \texttt{gate\_record}, \texttt{inventory},
\texttt{attestation}, \texttt{report}, \texttt{evidence\_package},
\texttt{operational\_record}, \texttt{build\_definition},
\texttt{release\_artifact}, \texttt{model}, \texttt{validation\_output},
\texttt{binding\_record}.

\textbf{EvidencePattern.} An expected evidence shape for deterministic
review support. Required fields: \texttt{evidence\_pattern\_id},
\texttt{canonical\_evidence\_kind}, \texttt{expectation},
\texttt{validation\_method}. EvidencePattern is defined and participates
in the artifact traceability chain
(\texttt{artifact\_supports\_evidence\_pattern}) but is not instantiated
as a first-class retrieval type in AppSec Core v0. It is instantiated in
the runtime layer described in companion work.

\paragraph{Relations}\label{relations}

\textbf{Table 2.} Relation types in AppSec Core v0, stable across all
slices and referenced consistently across all three companion papers.
Prose shorthand (\emph{realizes}, \emph{implements}, \emph{evidences})
appears in narrative descriptions as aliases for these identifiers; the
canonical form is the relation name in the table below.

\footnotesize
\setlength{\tabcolsep}{3pt}
\begin{longtable}{@{}L{5.0cm}L{3.0cm}L{2.7cm}L{3.9cm}@{}}
\toprule
\textbf{Relation} & \textbf{From} & \textbf{To} & \textbf{Meaning} \\
\midrule
\endfirsthead
\toprule
\textbf{Relation} & \textbf{From} & \textbf{To} & \textbf{Meaning} \\
\midrule
\endhead
\bottomrule
\endfoot
\bottomrule
\endlastfoot
\codeid{objective_realized_}\newline\codeid{by_practice} & ControlObjective & Practice & The practice operationalizes the objective \\
\codeid{objective_implemented_}\newline\codeid{by_mechanism} & ControlObjective & Mechanism & The mechanism technically implements the objective \\
\codeid{objective_expects_}\newline\codeid{artifact} & ControlObjective & Artifact & The artifact evidences achievement of the objective \\
\codeid{objective_verified_by_}\newline\codeid{evidence_pattern} & ControlObjective & EvidencePattern & The pattern defines how to verify the objective \\
\codeid{artifact_supports_}\newline\codeid{evidence_pattern} & Artifact & EvidencePattern & The artifact feeds the evidence pattern \\
\codeid{manual_requirement_}\newline\codeid{maps_to_objective} & ExternalRequirement & ControlObjective & Cross-framework mapping relation \\
\end{longtable}
\normalsize
\setlength{\tabcolsep}{4pt}

\paragraph{Traceability Chain}\label{traceability-chain}

The relations form a traceability chain from normative to evidentiary
(using active voice; the relation table above uses the equivalent
passive form). A \texttt{ControlObjective} is realized through
\texttt{Practice}, implemented through \texttt{Mechanism}, and evidenced
through \texttt{Artifact}; the associated \texttt{EvidencePattern}
specifies how that evidence is to be verified.

This chain ensures that each security concern is represented from
normative objective, through operational realization and technical
implementation, to evidentiary verification.

\paragraph{Instance Identifiers}\label{instance-identifiers}

Instance IDs follow a consistent prefix convention: \texttt{ACO-}
(ControlObjective), \texttt{ACP-} (Practice), \texttt{ACM-} (Mechanism),
\texttt{ACA-} (Artifact), followed by the slice abbreviation and a
sequence number. For example, \texttt{ACO-IVF-003} is the third control
objective in the Input Validation slice; \texttt{ACP-IVF-003} is the
corresponding practice.

\paragraph{Instance Example}\label{instance-example}

The following excerpt from slice ACO-IVF (Input Validation) illustrates
the four types and their relations:

\begin{Shaded}
\begin{Highlighting}[]
\FunctionTok{ControlObjective}\KeywordTok{:}
\AttributeTok{  }\FunctionTok{id}\KeywordTok{:}\AttributeTok{ ACO{-}IVF{-}003}
\AttributeTok{  }\FunctionTok{name}\KeywordTok{:}\AttributeTok{ Schema Validation Before Deserialization}
\AttributeTok{  }\FunctionTok{objective\_kind}\KeywordTok{:}\AttributeTok{ atomic}
\AttributeTok{  }\FunctionTok{objective\_type}\KeywordTok{:}\AttributeTok{ preventive}
\AttributeTok{  }\FunctionTok{statement}\KeywordTok{:}\AttributeTok{ }\StringTok{"All structured input must be validated against an explicit schema}
\StringTok{              before deserialization or internal processing"}
\AttributeTok{  }\FunctionTok{verification\_posture}\KeywordTok{:}\AttributeTok{ }\StringTok{"Verify typed model or JSON Schema at API entry point"}

\FunctionTok{Practice}\KeywordTok{:}
\AttributeTok{  }\FunctionTok{id}\KeywordTok{:}\AttributeTok{ ACP{-}IVF{-}003}
\AttributeTok{  }\FunctionTok{name}\KeywordTok{:}\AttributeTok{ Typed Model Validation At Entry Point}
\AttributeTok{  }\FunctionTok{practice\_family}\KeywordTok{:}\AttributeTok{ validation\_and\_analysis}
\AttributeTok{  }\FunctionTok{supports\_objectives}\KeywordTok{:}\AttributeTok{ }\KeywordTok{[}\AttributeTok{ACO{-}IVF}\DecValTok{{-}003}\KeywordTok{]}

\FunctionTok{Mechanism}\KeywordTok{:}
\AttributeTok{  }\FunctionTok{id}\KeywordTok{:}\AttributeTok{ ACM{-}IVF{-}003}
\AttributeTok{  }\FunctionTok{name}\KeywordTok{:}\AttributeTok{ Strict Typed Model Enforcement}
\AttributeTok{  }\FunctionTok{mechanism\_family}\KeywordTok{:}\AttributeTok{ validation\_and\_analysis}
\AttributeTok{  }\FunctionTok{supports\_practices}\KeywordTok{:}\AttributeTok{ }\KeywordTok{[}\AttributeTok{ACP{-}IVF}\DecValTok{{-}003}\KeywordTok{]}
\AttributeTok{  }\FunctionTok{note}\KeywordTok{:}\AttributeTok{ }\StringTok{"e.g., pydantic BaseModel with strict=True"}

\FunctionTok{Artifact}\KeywordTok{:}
\AttributeTok{  }\FunctionTok{id}\KeywordTok{:}\AttributeTok{ ACA{-}IVF{-}003}
\AttributeTok{  }\FunctionTok{name}\KeywordTok{:}\AttributeTok{ Schema Validation Test Suite}
\AttributeTok{  }\FunctionTok{canonical\_role}\KeywordTok{:}\AttributeTok{ validation\_output}
\AttributeTok{  }\FunctionTok{supports\_objectives}\KeywordTok{:}\AttributeTok{ }\KeywordTok{[}\AttributeTok{ACO{-}IVF}\DecValTok{{-}003}\KeywordTok{]}
\AttributeTok{  }\FunctionTok{note}\KeywordTok{:}\AttributeTok{ }\StringTok{"Test suite with malformed payload cases and schema coverage"}
\end{Highlighting}
\end{Shaded}

\subsubsection{3.5 Slice Contracts}\label{slice-contracts}

Each slice is governed by a formal contract. A contract specifies:

\begin{itemize}
\tightlist
\item
  \textbf{Scope}: what security concerns the slice covers
\item
  \textbf{Non-goals}: what is explicitly excluded (e.g., ASC-06 covers
  secret handling but not full cryptography; ASC-07 covers input
  handling but not runtime observability)
\item
  \textbf{Objective set}: the 7 control objectives with their
  \texttt{objective\_kind} and composite rules
\item
  \textbf{Module position}: which entity types are first-class (CO,
  Practice, Mechanism, Artifact, EvidencePattern), which are deferred
  (Threat, Signal)
\item
  \textbf{External alignment status}: explicitly deferred in v0
\end{itemize}

Contracts serve two functions: they \textbf{prevent scope creep} within
existing slices, and they \textbf{enable principled extension} ---
adding a new slice requires defining a contract, not modifying existing
ones.

\subsubsection{3.6 Cross-Slice Vocabulary}\label{cross-slice-vocabulary}

A shared vocabulary applies uniformly across all slices:

\begin{itemize}
\tightlist
\item
  \textbf{6 practice families} (governance\_and\_review,
  policy\_and\_gate\_enforcement, identity\_and\_trust\_control,
  validation\_and\_analysis, integrity\_traceability\_and\_records,
  rollout\_recovery\_and\_runtime\_records)
\item
  \textbf{15 canonical artifact roles} (configuration,
  governance\_record, review\_record, approval\_record, gate\_record,
  inventory, attestation, report, evidence\_package,
  operational\_record, build\_definition, release\_artifact, model,
  validation\_output, binding\_record)
\item
  \textbf{6 stable relations} connecting entity types
\end{itemize}

This shared vocabulary enables cross-slice queries and uniform analysis.
A tool that queries practices in one slice uses the same vocabulary and
structure in all slices.

\begin{center}\rule{0.5\linewidth}{0.5pt}\end{center}

\subsection{4. Normalization Approach}\label{normalization-approach}

\subsubsection{4.1 The Normalization
Problem}\label{the-normalization-problem}

Given two external requirements \emph{r\_1} ∈ Framework A and
\emph{r\_2} ∈ Framework B, the normalization question is: \textbf{do
\emph{r\_1} and \emph{r\_2} address the same security concern?} Without
a shared semantic model, this question can only be answered by reading
both requirements and making a human judgment --- a process that does
not scale.

AppSec Core provides a principled answer: \emph{r\_1} and \emph{r\_2}
address the same concern if they map to the same AppSec Core control
objective(s). The ontology is the \textbf{tertium comparationis} --- the
shared reference point that makes cross-framework comparison possible.

\subsubsection{4.2 Mapping Protocol}\label{mapping-protocol}

For each external framework, the mapping proceeds as:

\begin{enumerate}
\def\labelenumi{\arabic{enumi}.}
\tightlist
\item
  \textbf{Slice assignment}: each external requirement is assigned to
  one or more AppSec Core slices based on semantic scope
\item
  \textbf{Objective mapping}: within each slice, the requirement is
  mapped to specific control objectives with an explicit strength
  indicator:

  \begin{itemize}
  \tightlist
  \item
    \texttt{primary}: the objective is the main semantic anchor for this
    requirement
  \item
    \texttt{secondary}: the objective is materially relevant but not the
    best primary anchor
  \end{itemize}
\item
  \textbf{Rationale recording}: each mapping includes a human-authored
  rationale explaining the semantic link
\end{enumerate}

Mappings are \texttt{many-to-one} and \texttt{one-to-many}: a single
framework requirement may map to multiple objectives, and a single
objective may be the target of multiple framework requirements. This is
expected --- normalization is not a bijection.

In this study, mappings were performed by the ontology's authors (a
limitation we acknowledge in Section 8). The unit of analysis varied by
framework: SSDF practices (N=20), ASVS thematic clusters (N=36), SLSA
requirements (N=14), CIS controls (N=8 of 18), CAPEC patterns (N=13 of
View 683). For CIS and CAPEC, whose scope extends beyond application
security, we applied an explicit inclusion rule: an item was included if
it directly affects how software is specified, developed, built, tested,
deployed, or monitored. Items addressing enterprise IT administration,
physical security, network infrastructure, or non-software-operational
concerns were excluded. This rule yielded 8 CIS controls (e.g., CIS-4
Secure Configuration, CIS-16 Application Software Security) and 13 CAPEC
patterns (the supply chain attack view). SSDF, ASVS, and SLSA were
included in full, as their scope is inherently within the AppSec
boundary.

A requirement was considered to match an objective if it addressed the
same security concern at a comparable level of abstraction (scope
overlap) and the mapping was operationally meaningful --- not merely
topically adjacent.

\subsubsection{4.3 Semantic Convergence}\label{semantic-convergence}

When requirements from different frameworks map to the same objective,
we observe \textbf{semantic convergence}: distinct vocabularies
describing the same underlying concern. This convergence is the
empirical signature of normalization.

\subsubsection{4.4 The Normalization
Property}\label{the-normalization-property}

We define normalization precisely:

\begin{quote}
Two external requirements \emph{r₁ ∈ F\_A} and \emph{r₂ ∈ F\_B} are
\textbf{normalized-equivalent with respect to a concern} if they share
at least one primary ControlObjective mapping within the same AppSec
Core slice. They may additionally overlap on secondary objectives; the
primary mapping anchors the normalized concern.
\end{quote}

Once mapped, downstream reasoning operates on the \textbf{normalized
objective layer}, not on framework-specific vocabulary:

\begin{itemize}
\tightlist
\item
  \textbf{coverage analysis} --- coverage is assessed against the
  objective set; a requirement is covered when its primary objective(s)
  are satisfied
\item
  \textbf{implementation reasoning} --- the practice and mechanism
  chains associated with the objective apply regardless of which
  framework motivated the requirement
\item
  \textbf{verification reasoning} --- the evidence patterns and
  artifacts are determined by the objective, not by the originating
  framework
\end{itemize}

The canonical traceability structure is as follows: an
\texttt{ExternalRequirement} maps to a \texttt{ControlObjective}; that
objective is operationalized through \texttt{Practice} and
\texttt{Mechanism}, and it is evidenced through \texttt{Artifact} and
\texttt{EvidencePattern}.

This is the property that distinguishes normalization from
classification. A taxonomy assigns items to categories but preserves
framework-specific identity. Normalization performs \textbf{canonical
reduction}: requirements addressing the same security concern are
reduced to a shared operational representation anchored on common
objectives, and subsequent coverage analysis operates on that
representation. The framework origin becomes provenance metadata, not an
input to reasoning.

The reduction is not strict set equality --- different frameworks may
cover different subsets of a slice's objectives (as shown in Section
5.1). What is shared is the \textbf{normalized concern}: the primary
objective(s) and their associated operational chains. This is sufficient
for unified coverage analysis while honestly reflecting the granularity
differences between frameworks.

\begin{center}\rule{0.5\linewidth}{0.5pt}\end{center}

\subsection{5. Mapping Examples}\label{mapping-examples}

\subsubsection{5.1 Configuration Security: ASVS, SSDF, and
CIS}\label{configuration-security-asvs-ssdf-and-cis}

Three frameworks address configuration and secret security:

\begin{itemize}
\tightlist
\item
  \textbf{ASVS v5.0}: cluster \texttt{secure\_configuration\_baseline}
  --- verification of secure defaults, secret externalization
\item
  \textbf{SSDF PW.9}: \emph{``Configure Software to Have Secure Settings
  by Default''}
\item
  \textbf{CIS-4}: \emph{``Secure Configuration of Enterprise Assets and
  Software''}
\end{itemize}

All three converge on slice \textbf{ACO-SPC} (Secret Handling, Protected
Configuration \& Operational Identities):

\begin{longtable}{@{}L{2.25cm}L{7.2cm}C{1cm}C{1cm}C{1cm}@{}}
\toprule
\textbf{Objective} & \textbf{Description} & \textbf{ASVS} & \textbf{SSDF} & \textbf{CIS} \\
\midrule
\endfirsthead
\toprule
\textbf{Objective} & \textbf{Description} & \textbf{ASVS} & \textbf{SSDF} & \textbf{CIS} \\
\midrule
\endhead
\bottomrule
\endfoot
\bottomrule
\endlastfoot
ACO-SPC-001 & Secret leak prevention (anti-hardcoding, scanning) & $\checkmark$ & $\checkmark$ & $\checkmark$ \\
ACO-SPC-002 & Vault-backed protected storage and controlled retrieval & $\checkmark$ & $\checkmark$ & --- \\
ACO-SPC-003 & Rotation and renewal discipline & $\checkmark$ & --- & --- \\
ACO-SPC-004 & Operational identity binding (OIDC, short-lived credentials) & --- & --- & $\checkmark$ \\
\end{longtable}

The convergence is not uniform: ASVS provides the most granular
decomposition; SSDF operates at practice level; CIS includes enterprise
concerns beyond AppSec scope. Yet all three reduce to the same
normalized concern within ACO-SPC, anchored by overlapping
ControlObjectives.

The operational representation is shared:

\begin{Shaded}
\begin{Highlighting}[]
\FunctionTok{ACP{-}SPC{-}001}\KeywordTok{:}\CommentTok{  \# Practice}
\AttributeTok{  }\FunctionTok{name}\KeywordTok{:}\AttributeTok{ Secret Leak Prevention In Source And Pipeline}
\AttributeTok{  }\FunctionTok{supports\_objectives}\KeywordTok{:}\AttributeTok{ }\KeywordTok{[}\AttributeTok{ACO{-}SPC}\DecValTok{{-}001}\KeywordTok{]}

\FunctionTok{ACP{-}SPC{-}002}\KeywordTok{:}\CommentTok{  \# Practice}
\AttributeTok{  }\FunctionTok{name}\KeywordTok{:}\AttributeTok{ Vault{-}Backed Secret Storage}
\AttributeTok{  }\FunctionTok{supports\_objectives}\KeywordTok{:}\AttributeTok{ }\KeywordTok{[}\AttributeTok{ACO{-}SPC}\DecValTok{{-}002}\KeywordTok{]}
\end{Highlighting}
\end{Shaded}

Requirements from the manual's \texttt{requirements\_catalog} (CFG-003:
no hardcoded parameters; CFG-004: externalized configuration; ENC-006:
secret exposure detection; ENC-007: key rotation) also map to these same
objectives --- demonstrating that practitioner-derived requirements and
framework requirements converge on the same representation.

This example illustrates the normalization property (Section 4.4): ASVS,
SSDF, and CIS requirements share primary objectives within ACO-SPC
(notably ACO-SPC-001), providing a shared operational basis for the
normalized concern of configuration security. Coverage is assessed once
at the objective level: an organization that satisfies the mapped
objectives addresses the underlying concern expressed by all three
frameworks --- subject to the correctness and granularity of the
mapping. After the mapping step, downstream coverage analysis operates
on the normalized objective layer, not on three separate
framework-specific vocabularies.

\subsubsection{5.2 Input Validation: ASVS, SSDF, and
CAPEC}\label{input-validation-asvs-ssdf-and-capec}

\begin{itemize}
\tightlist
\item
  \textbf{ASVS v5.0}: clusters \texttt{injection\_and\_sanitization},
  \texttt{input\_contract\_validation},
  \texttt{validation\_before\_internal\_use}
\item
  \textbf{SSDF PW.5}: \emph{``Create Source Code by Adhering to Secure
  Coding Practices''}
\item
  \textbf{CAPEC-536}: \emph{``Data Injected During Configuration''}
\end{itemize}

All three converge on slice \textbf{ACO-IVF} (Input Validation, Safe
Parsing \& Controlled Failure):

\begin{longtable}{@{}L{2.25cm}L{7.2cm}C{1cm}C{1cm}C{1cm}@{}}
\toprule
\textbf{Objective} & \textbf{Description} & \textbf{ASVS} & \textbf{SSDF} & \textbf{CAPEC} \\
\midrule
\endfirsthead
\toprule
\textbf{Objective} & \textbf{Description} & \textbf{ASVS} & \textbf{SSDF} & \textbf{CAPEC} \\
\midrule
\endhead
\bottomrule
\endfoot
\bottomrule
\endlastfoot
ACO-IVF-001 & Entry-point input contract validation & $\checkmark$ & $\checkmark$ & --- \\
ACO-IVF-002 & Schema, type, and allowlist discipline & $\checkmark$ & $\checkmark$ & --- \\
ACO-IVF-003 & Dangerous-pattern exclusion & $\checkmark$ & $\checkmark$ & $\checkmark$ \\
ACO-IVF-004 & Validation before internal use & $\checkmark$ & --- & $\checkmark$ \\
\end{longtable}

ASVS provides the most granular decomposition. SSDF provides a broader
practice-level mapping. CAPEC provides an attack-pattern perspective on
the same concerns. AppSec Core absorbs all three into the same
normalized concern, represented by overlapping objectives within the
slice.

\subsubsection{5.3 Supply Chain: SLSA, SSDF, and
CAPEC}\label{supply-chain-slsa-ssdf-and-capec}

\begin{itemize}
\tightlist
\item
  \textbf{SLSA Build Track}: provenance, build integrity, platform
  isolation
\item
  \textbf{SSDF PW.3 / PS.2}: verify third-party software, verify release
  integrity
\item
  \textbf{CAPEC-185/446/206}: malicious download, malicious component,
  signing attack
\end{itemize}

All three converge on slice \textbf{ACO-SCBI} (Supply Chain \& Build
Integrity):

\begin{longtable}{@{}L{2.5cm}L{6.95cm}C{0.9cm}C{0.9cm}C{1cm}@{}}
\toprule
\textbf{Objective} & \textbf{Description} & \textbf{SLSA} & \textbf{SSDF} & \textbf{CAPEC} \\
\midrule
\endfirsthead
\toprule
\textbf{Objective} & \textbf{Description} & \textbf{SLSA} & \textbf{SSDF} & \textbf{CAPEC} \\
\midrule
\endhead
\bottomrule
\endfoot
\bottomrule
\endlastfoot
ACO-SCBI-001 & Dependency inventory and SBOM traceability & $\checkmark$ & $\checkmark$ & --- \\
ACO-SCBI-002 & Dependency risk analysis and policy gating & --- & $\checkmark$ & $\checkmark$ \\
ACO-SCBI-003 & Approved source and registry governance & $\checkmark$ & --- & $\checkmark$ \\
ACO-SCBI-004 & Trusted build definition and execution integrity & $\checkmark$ & --- & --- \\
ACO-SCBI-005 & Promotion gates and approval governance & $\checkmark$ & $\checkmark$ & --- \\
ACO-SCBI-006 & Artifact provenance and signature integrity & $\checkmark$ & $\checkmark$ & $\checkmark$ \\
\end{longtable}

This is the most convergent slice: all three frameworks contribute to 4+
of 6 atomic objectives, despite originating from different domains
(build integrity, development practices, attack patterns).

\subsubsection{5.4 Convergence Metrics}\label{convergence-metrics}

Table 3 aggregates the convergence data from the three presented slices.

\textbf{Table 3.} Cross-framework convergence per slice (based on
presented mapping tables).

\begin{longtable}{@{}L{2.2cm}C{2.3cm}C{2.8cm}C{2.8cm}C{2cm}@{}}
\toprule
\textbf{Slice} & \textbf{Objectives shown} & \textbf{$\geq$ 2 frameworks} & \textbf{$\geq$ 3 frameworks} & \textbf{Max overlap} \\
\midrule
\endfirsthead
\toprule
\textbf{Slice} & \textbf{Objectives shown} & \textbf{$\geq$ 2 frameworks} & \textbf{$\geq$ 3 frameworks} & \textbf{Max overlap} \\
\midrule
\endhead
\bottomrule
\endfoot
\bottomrule
\endlastfoot
ACO-SPC & 4 & 2 (50\%) & 1 (25\%) & 3 \\
ACO-IVF & 4 & 4 (100\%) & 1 (25\%) & 3 \\
ACO-SCBI & 6 & 5 (83\%) & 1 (17\%) & 3 \\
\textbf{Total} & \textbf{14} & \textbf{11 (79\%)} & \textbf{3 (21\%)} & \textbf{3} \\
\end{longtable}

Across the 14 objectives analyzed, 79\% are mapped by at least two
frameworks, and each slice contains at least one objective mapped by all
three frameworks in its group. The triple-convergence points ---
ACO-SPC-001 (secret leak prevention), ACO-IVF-003 (dangerous-pattern
exclusion), ACO-SCBI-006 (artifact provenance) --- represent core
concerns where framework agreement is strongest and where canonical
reduction provides the greatest value.

Single-framework objectives (3 of 14, 21\%) occur at specialization
edges: rotation discipline (ASVS only), operational identity binding
(CIS only), trusted build definition (SLSA only). The ontology
accommodates these framework-specific extensions alongside the shared
core --- it does not force convergence where the domain is genuinely
specialized.

We note that these three slices were selected for highest
cross-framework overlap; convergence rates for the remaining seven
slices may be lower. The metrics above should be read as evidence of
convergence on core concerns, not as a claim about the full ontology.

\subsubsection{5.5 Coverage Summary}\label{coverage-summary}

We present mapping examples for three slices (ACO-SPC, ACO-IVF,
ACO-SCBI). The remaining seven slices have been validated through
internal mapping exercises but are not presented here for space reasons.

Table 4 summarizes framework coverage across the full mapping exercise.

\textbf{Table 4.} Framework coverage summary.

\begin{longtable}{@{}L{3.55cm}L{3.15cm}C{2cm}C{1.8cm}C{2.05cm}@{}}
\toprule
\textbf{Framework} & \textbf{Unit of analysis} & \textbf{Items mapped} & \textbf{Slices used} & \textbf{Items outside model} \\
\midrule
\endfirsthead
\toprule
\textbf{Framework} & \textbf{Unit of analysis} & \textbf{Items mapped} & \textbf{Slices used} & \textbf{Items outside model} \\
\midrule
\endhead
\bottomrule
\endfoot
\bottomrule
\endlastfoot
SSDF v1.1 & Practices & 20 & 8 & 0 \\
ASVS v5.0 & Thematic clusters & 36 & 10 & 0 \\
SLSA v1.0 & Requirements & 14 & 3 & 0 \\
CIS v8.1 (AppSec scope) & Controls & 8 & 6 & 0 \\
CAPEC v3.9 (View 683) & Attack patterns & 13 & 4 & 0 \\
\textbf{Total} &  & \textbf{91} & \textbf{10} & \textbf{0} \\
\end{longtable}

Every external requirement mapped to at least one AppSec Core slice. No
requirement fell outside the 10-slice model under the stated inclusion
rules and mapping protocol (Section 4.2). This provides evidence --- not
proof --- of \emph{framework coverage}: the model accommodates these
frameworks' requirements.

We distinguish this from \emph{domain coverage} --- whether AppSec Core
captures all relevant AppSec concerns. The latter would require
validation against a broader set of practitioners, organizations, and
domains (e.g., embedded systems, ML pipeline security), which is future
work. The evidence here supports framework coverage for the analyzed
sources; domain coverage remains an open question.

\begin{center}\rule{0.5\linewidth}{0.5pt}\end{center}

\subsection{6. Structural Properties}\label{structural-properties}

\subsubsection{6.1 Structural Invariance}\label{structural-invariance}

Every slice exhibits the same structure: 7 objectives, practices linked
via \texttt{objective\_realized\_by\_practice}, mechanisms linked via
\texttt{objective\_implemented\_by\_mechanism}, artifacts linked via
\texttt{objective\_expects\_artifact}, and a governing contract. This
invariance has practical consequences:

\begin{itemize}
\tightlist
\item
  \textbf{Tooling reuse}: a tool that analyzes one slice works on all
  slices unchanged
\item
  \textbf{Uniform coverage analysis}: the same algorithm determines
  coverage for any domain
\item
  \textbf{Predictable extension}: new slices follow the same template
\end{itemize}

\subsubsection{6.2 Framework Independence}\label{framework-independence}

AppSec Core was derived from practitioner semantics. This means:

\begin{itemize}
\tightlist
\item
  Adding a new framework (e.g., ASVS v6.0) requires mapping its
  requirements to existing objectives, not restructuring slices
\item
  The 10-slice model reflects what practitioners do; framework
  vocabulary is absorbed, not adopted
\item
  Security concerns evolve slowly; framework vocabularies evolve
  rapidly. The ontology models the former.
\end{itemize}

\subsubsection{6.3 Extensibility}\label{extensibility}

The contract-based model enables principled extension. New slices can be
defined when the domain requires it --- the current count of 10 is
methodological, not ontological. Existing contracts prevent scope creep.
The cross-slice vocabulary ensures new slices integrate cleanly.

\subsubsection{6.4 What AppSec Core is
Not}\label{what-appsec-core-is-not}

\begin{itemize}
\tightlist
\item
  \textbf{Not a replacement for frameworks}: organizations still use
  SSDF, ASVS, etc. for their specific purposes. AppSec Core provides the
  normalization layer for cross-framework analysis.
\item
  \textbf{Not a maturity model}: it does not assess how well practices
  are performed.
\item
  \textbf{Not a vulnerability taxonomy}: it models the solution space,
  not the problem space.
\item
  \textbf{Not a formal ontology} in the OWL/description-logic sense
  {[}16, 17, 19{]}: it is a typed semantic model with explicit
  relations, optimized for practitioner use and tooling integration. The
  informal use of \emph{ontology} follows Gruber's sense of ``a
  specification of a conceptualization'' {[}16{]}.
\end{itemize}

\begin{center}\rule{0.5\linewidth}{0.5pt}\end{center}

\subsection{7. Discussion}\label{discussion}

\subsubsection{7.1 Normalization
vs.~Classification}\label{normalization-vs.-classification}

The distinction is precise:

\begin{itemize}
\item
  \textbf{Classification} assigns requirements to categories. Each
  requirement retains its framework-specific identity; the categories
  provide grouping but not reduction. Analysis still requires reading
  each framework's vocabulary.
\item
  \textbf{Normalization} performs \textbf{canonical reduction}:
  requirements from different frameworks that address the same security
  concern are reduced to shared ControlObjectives and their associated
  operational chains --- shared practices, shared mechanisms, shared
  evidence --- such that they become operationally equivalent for
  coverage analysis. The framework of origin becomes provenance
  metadata, not an input to downstream reasoning.
\end{itemize}

The evidence in Section 5 supports that AppSec Core performs
normalization, not merely classification. SSDF PW.9, ASVS
\texttt{secure\_configuration\_baseline}, and CIS-4 are three distinct
vocabulary entries that, after mapping, resolve to the same normalized
concern within ACO-SPC, anchored by overlapping ControlObjectives and
shared operational chains. Downstream coverage analysis operates on the
canonical chain --- not on three separate framework-specific
assessments. This is canonical reduction, not grouping.

\subsubsection{7.2 Towards
Standardization}\label{towards-standardization}

AppSec Core is proposed as a \textbf{candidate} for standardization, not
as a finished standard. The path would require:

\begin{enumerate}
\def\labelenumi{\arabic{enumi}.}
\tightlist
\item
  \textbf{Multi-organization derivation}: the current ontology reflects
  one team's experience. Cross-organizational validation would test
  whether the 10-slice model generalizes.
\item
  \textbf{Requirement-level alignment}: v0 maps at the practice/cluster
  level. A standard would require finer-grained alignment.
\item
  \textbf{Community governance}: slice contracts are designed to support
  community-driven extension, but this has not been tested.
\item
  \textbf{Formal specification}: the current YAML representation may
  need formalization (e.g., SHACL {[}18{]}, OWL {[}17{]}) for
  interoperability, though accessibility must be weighed against
  formality.
\end{enumerate}

\subsubsection{7.3 Practical Benefits}\label{practical-benefits}

If AppSec Core or a similar model were adopted:

\begin{itemize}
\tightlist
\item
  \textbf{Unified compliance surface}: one coverage analysis for N
  frameworks
\item
  \textbf{Version resilience}: framework updates require re-mapping, not
  restructuring
\item
  \textbf{Tool interoperability}: security tools could align on shared
  identifiers
\item
  \textbf{Cross-organization comparison}: organizations using the same
  normalization layer could compare coverage without exposing
  proprietary documentation
\end{itemize}

\subsubsection{7.4 The Practitioner-Derived
Argument}\label{the-practitioner-derived-argument}

A normalization ontology must be independent of the frameworks it
normalizes; otherwise it is a biased projection. AppSec Core avoids this
by being derived from what practitioners actually do --- documented
before many of the frameworks it now normalizes.

The SbD-ToE corpus draws on a prescriptive Security by Design framework
first documented in practitioner deployment in 2018 {[}19{]}, predating
the publication of SSDF v1.0 (2021), SLSA v1.0 (2021), and CIS Controls
v8 (2021). Several concerns now formalised in those frameworks ---
including risk-proportional requirement selection and multi-framework
normalisation into a unified prescriptive activity model --- were
already operationalised and documented in client-facing practice from
that period. This does not predate the underlying security domain
knowledge; it predates the formal framework publications. The frameworks
codified concerns that practitioners had already identified and
implemented.

More recent examples corroborate the same pattern at the
individual-concept level: secure configuration baselines (CFG-001→007)
were documented in the corpus in 2023; ASVS v5.0 formalised
\texttt{secure\_configuration\_baseline} in 2025. Container admission
policies were implemented in practice in 2022; SLSA formalised build
platform isolation in 2023.

This temporal precedence supports the independence claim: AppSec Core
was not reverse-engineered from frameworks to fit their structure. It
was abstracted from operational practice, and the frameworks
independently converged on the same domain concerns.

\begin{center}\rule{0.5\linewidth}{0.5pt}\end{center}

\subsection{8. Limitations}\label{limitations}

\textbf{Single-team derivation.} The ontology reflects one team's 10+
years of experience. The 10-slice structure may not capture concerns
unique to domains not encountered by this team (e.g., embedded systems,
ML pipeline security, mobile-native development). We argue for
structural transferability --- the typed schema is domain-independent;
the slices may need adjustment --- but this has not been validated.

\textbf{V0 granularity.} External alignment operates at the
practice/cluster level, not at individual requirement level. The
convergence evidence in Section 5 uses this granularity.

\textbf{No formal specification.} The ontology is represented in YAML
with explicit types and relations. This is a deliberate trade-off
between practitioner accessibility and formal interoperability. A formal
OWL or SHACL representation is future work.

\textbf{Evaluator subjectivity.} The mapping exercise was performed by
the ontology's authors. Inter-rater reliability with independent
evaluators was not measured.

\textbf{Temporal specificity.} The framework landscape is captured as of
early 2026.

\begin{center}\rule{0.5\linewidth}{0.5pt}\end{center}

\subsection{9. Conclusion}\label{conclusion}

We have presented AppSec Core, a normalized ontology for application
security with 10 domain slices and 234 typed instances. The ontology
provides:

\begin{itemize}
\tightlist
\item
  \textbf{Measured cross-framework convergence}: in the presented
  slices, 79\% of analyzed ControlObjectives are mapped by at least two
  frameworks, supporting the claim that core AppSec concerns can be
  reduced into shared objectives despite vocabulary differences
\item
  \textbf{Structural invariance}: every slice uses the same four-type
  schema and relation model
\item
  \textbf{Framework independence}: derived from practitioner semantics,
  stable as frameworks evolve
\item
  \textbf{Contract-based governance}: each slice governed by a formal
  contract supporting extension
\end{itemize}

The evidence supports that AppSec Core functions as a normalization
layer, not merely a taxonomy. It provides a \textbf{canonical reduction
layer}: heterogeneous framework requirements addressing the same
security concerns are reduced to shared ControlObjectives and their
associated operational chains. Downstream coverage analysis operates on
this normalized objective layer independently of framework vocabulary.
This is canonical reduction --- the defining property of normalization
--- not grouping or alignment. The reduction does not claim full
normative equivalence between frameworks; it claims a shared operational
basis sufficient for unified coverage analysis.

To our knowledge, the application security domain lacks a shared
semantic model for cross-framework normalization. AppSec Core is our
proposal for this model. It is not a replacement for existing frameworks
--- it is a candidate for the canonical layer between them.

\begin{center}\rule{0.5\linewidth}{0.5pt}\end{center}

\subsection{10. Artifact Availability}\label{artifact-availability}

Curated supporting artifacts for this paper are available in the
companion public repository at
\url{https://github.com/sbd-ai-runtime/appsec-core-ontology-research}.
For this paper, the relevant materials are organized under
\texttt{papers/01-appsec-core-normalized-ontology/artifacts/}, notably
\texttt{papers/01-appsec-core-normalized-ontology/artifacts/ontology/},
\texttt{papers/01-appsec-core-normalized-ontology/artifacts/schema/},
and
\texttt{papers/01-appsec-core-normalized-ontology/artifacts/slice\_contracts/}.
The same repository also contains this paper's curated source under
\texttt{papers/01-appsec-core-normalized-ontology/source/} and its
public PDF under
\texttt{papers/01-appsec-core-normalized-ontology/pdf/}.

\begin{center}\rule{0.5\linewidth}{0.5pt}\end{center}

\subsection{References}\label{references}

{[}1{]} NIST. Secure Software Development Framework (SSDF) Version 1.1.
SP 800-218. 2022.

{[}2{]} OWASP. Application Security Verification Standard (ASVS) v5.0.0.
2025.

{[}3{]} SLSA. Supply-chain Levels for Software Artifacts. Build Track
v1.0. 2023.

{[}4{]} CIS. Center for Internet Security Controls v8.1. 2024.

{[}5{]} European Union. Digital Operational Resilience Act (DORA).
Regulation (EU) 2022/2554. 2022.

{[}6{]} European Union. Network and Information Security Directive
(NIS2). Directive (EU) 2022/2555. 2022.

{[}7{]} A. Herzog, N. Shahmehri, and C. Duma, ``An ontology of
information security,'' \emph{Int. J. Inf. Security and Privacy},
vol.~1, no. 4, pp.~1--23, 2007.

{[}8{]} C. Blanco, J. Lasheras, R. Valencia-García, E. Fernández-Medina,
A. Toval, and M. Piattini, ``A systematic review and comparison of
security ontologies,'' in \emph{Proc. 3rd Int. Conf. Availability,
Reliability and Security (ARES)}, 2008, pp.~813--820, doi:
10.1109/ARES.2008.33.

{[}9{]} MITRE. Common Weakness Enumeration (CWE). Available:
https://cwe.mitre.org (accessed: 2026-04-06).

{[}10{]} MITRE. Common Attack Pattern Enumeration and Classification
(CAPEC). Available: https://capec.mitre.org (accessed: 2026-04-06).

{[}11{]} OWASP. Software Assurance Maturity Model (SAMM) v2.1. 2024.

{[}12{]} Synopsys. Building Security In Maturity Model (BSIMM) 13. 2023.

{[}13{]} NIST. Security and Privacy Controls for Information Systems and
Organizations. SP 800-53 Rev.~5. 2020.

{[}14{]} B. Fabian, S. Gürses, M. Heisel, T. Santen, and H. Schmidt, ``A
comparison of security requirements engineering methods,''
\emph{Requirements Engineering}, vol.~15, no. 1, pp.~7--40, 2010, doi:
10.1007/s00766-009-0092-x.

{[}15{]} A. Goknil, I. Kurtev, K. van den Berg, and J.-W. Veldhuis,
``Semantics of trace relations in requirements models for consistency
checking and inferencing,'' \emph{Software \& Systems Modeling},
vol.~10, no. 1, pp.~31--54, 2011, doi: 10.1007/s10270-009-0142-3.

{[}16{]} T. R. Gruber, ``A translation approach to portable ontology
specifications,'' \emph{Knowledge Acquisition}, vol.~5, no. 2,
pp.~199--220, 1993, doi: 10.1006/knac.1993.1008.

{[}17{]} W3C. OWL 2 Web Ontology Language: Document Overview (Second
Edition). W3C Recommendation, 2012. Available:
https://www.w3.org/TR/owl2-overview/ (accessed: 2026-04-06).

{[}18{]} W3C. Shapes Constraint Language (SHACL). W3C Recommendation,
2017. Available: https://www.w3.org/TR/shacl/ (accessed: 2026-04-06).

{[}19{]} P. Farinha, ``Framework Shiftleft --- Security by Design,''
Shiftleft internal technical documentation, 2018. Multiple client
deployment instances. Available from the authors upon request.

{[}20{]} F. Baader, I. Horrocks, and U. Sattler, ``An introduction to
description logics,'' in \emph{The Description Logic Handbook: Theory,
Implementation and Applications}, F. Baader, D. Calvanese, D.
McGuinness, D. Nardi, and P. F. Patel-Schneider, Eds. Cambridge
University Press, 2003, ch.~1, pp.~1--44.

\end{document}
